\documentclass[11pt,a4paper]{article}
\usepackage[T1]{fontenc}
\usepackage[utf8]{inputenc}
\usepackage[provide=*,french]{babel}
\usepackage{lmodern,microtype}
\usepackage[a4paper,margin=23mm,headheight=15pt]{geometry}
\usepackage{amsmath,amssymb,booktabs,tabularx,array,enumitem,fancyhdr,xcolor}
\usepackage[colorlinks=true,linkcolor=blue!40!black,urlcolor=blue!40!black]{hyperref}
\usepackage{xurl}
\input{production_metrics}
\definecolor{accent}{HTML}{173B5B}
\pagestyle{fancy}\fancyhf{}
\fancyhead[L]{\small\color{accent} P1 -- Note de révision V1.5}
\fancyhead[R]{\small 31 octobre 2026}
\fancyfoot[C]{\thepage}
\renewcommand{\headrulewidth}{0.3pt}
\setlength{\parindent}{0pt}
\setlength{\parskip}{6pt}
\setlength{\emergencystretch}{2em}
\setlist{itemsep=4pt,topsep=4pt,leftmargin=18pt}
\renewcommand{\arraystretch}{1.13}
\newcommand{\E}{\mathbb E}
\newcommand{\Law}{\mathcal L}
\newcommand{\eps}{\varepsilon}
\hypersetup{pdftitle={P1 V1.5 -- Note de revision et reponse aux retours},
pdfauthor={Brice Pouly},pdfsubject={Attributions, preuves developpees, marges et controle de la revision}}
\widowpenalty=10000\clubpenalty=10000
\begin{document}
{\Large\bfseries\color{accent} P1 V1.5 : réponse aux deux retours}\\[4pt]
{\large Note de révision -- 31 octobre 2026}

\section*{1. Ce que cette version change}
La V1.5 conserve l'organisation en \textbf{deux documents associés} :
l'article principal compte \textbf{\MainPageCount{} pages}, le compagnon
\textbf{\CompanionPageCount{} pages}. Elle précise les attributions,
développe un dénombrement conditionnel dense et rend explicites les
dépendances et les marges de la démonstration. Elle ne revendique pas
de nouvel exposant par rapport à la V1.4.

La révision répond à l'\emph{Audit bibliographique et de priorité de P1 V1.4}
et à la \emph{Contre-review de P1 V1.4}, tous deux datés du
31 octobre 2026. Ces retours ne fournissent pas de contre-exemple établi
ni de lacune bloquante. Leurs demandes ont été vérifiées dans les passages
concernés et, pour les nouvelles attributions, dans les textes primaires.
Ils restent des documents de travail éditorial et ne sont pas ajoutés à
la bibliographie scientifique.

\subsection*{Les résultats centraux conservés}
Avec \(\mu_N=\E Z_N\), \(\alpha_N=\operatorname{Var}(Z_N)-\mu_N\),
\(\Pi_\mu=\operatorname{Pois}(\mu)\), \(\Delta=\delta_1-\delta_0\) et
\[
 Q_{\mu,\alpha}=\Pi_\mu+\frac{\alpha}{2}\Pi_\mu*\Delta^{*2},
\]
le théorème 1.1 conserve, pour toute marge fixe \(\eps>0\),
\[
 c_{C,\eps}N^{-4/3-\eps}
 \le \|\Law(Z_N)-Q_{\mu_N,\alpha_N}\|_1
 \le C_{C,\eps}N^{-4/3+\eps}.
\]
Le troisième cumulant factoriel conserve la puissance optimale
\(41/30\). Pour les références explicites retenant les canaux jusqu'à
une hauteur \(h\) fixée, la puissance reste \(h/(h+1)\), avec
l'équivalent donné par le premier niveau omis.

Les \textbf{\NumberedResultCount{} énoncés numérotés} de type théorème,
proposition, lemme ou corollaire gardent leur document et leur numéro.
Les \textbf{\PreservedLabelCount{} repères de la V1.4} sont conservés.
La correspondance complète V1.4--V1.5 figure dans les sources.

\subsection*{Le sens des corrections}
Trois outils sont désormais attribués directement à leurs sources :
la borne binaire de Griesmer, l'expansion abstraite de Zeilberger et le
comptage quadratique uniforme de Cilleruelo--Garaev, également formulé
par Tao. Les comparaisons avec les approximations discrètes classiques
sont précisées. La contribution de P1 est ainsi décrite au niveau des
coefficients et configurations arithmétiques, des restes en norme et
des échelles optimales pour les références spécifiées.

\clearpage
\section*{2. Réponses aux demandes techniques}

\subsection*{R1. Développer le comptage à deux fenêtres}
\textbf{Traitement : intégré dans le lemme 15.2 et sa preuve en T11.}
Le résumé dans l'article est complété par une preuve structurée dans le
compagnon. Elle explicite les étapes suivantes.

Le noyau de l'application somme
\(W\oplus S_y\oplus S_z\to W+S_y+S_z\) se projette injectivement
sur chacune des deux cibles, puisque \(S_y\cap W=S_z\cap W=0\).
Après régularisation, les étiquettes de chaque cible engendrent le dual
du noyau restreint. Chaque composante binaire ne porte qu'une étiquette ;
le nombre de composantes rencontrant \emph{chaque} cible est donc minoré
séparément par la dimension restante. Les sommets actifs isolés sont inclus.

Les circuits de premiers rencontrant au moins trois blocs doivent toucher
une source fixée, car seules deux cibles sont libres. Les produits de
premiers correspondants divisent alors des valeurs de source fixées.
La preuve distingue le choix des paramètres qui bornent la hauteur des
coefficients du choix ultérieur du nombre d'équations géométriques.

Pour un alignement faible comportant \(u\) unités, l'inégalité
\(u-2\le\dim(S_y\cap S_z)\) contrôle le nombre de composantes à exclure :
au plus \(2(u-1)\). Il s'agit de retirer des composantes de la collection
sélectionnable, sans imposer une nouvelle restriction au noyau.
Si \(m_t\) composantes subsistent dans la cible \(t\), celles qui
n'utilisent qu'une coordonnée sont au nombre de
\[
 n_{1,t}\ge 2m_t-B,\qquad t=y,z,\quad B=L+1.
\]
Au-dessus du seuil \((1/2+\delta)L\), cette borne force des occurrences
à une seule coordonnée dans les deux cibles. Elle permet le dénombrement
renforcé, puis le passage aux poids \(p_yp_z2^d\).

\subsection*{R2. Rendre les petites pertes vérifiables}
\textbf{Traitement : registre des marges en section 1.5 et renvois dans
les sections 14--17.} La marge finale demandée est distinguée des
tolérances internes. Les écarts \(1/120\), \(1/60\), \(1/30\) et
\(2/15\) sont calculés explicitement ; la page suivante les rassemble.
Le registre de sources détaille également les quatre contributions de
la queue uniforme et les termes situés exactement à l'échelle \(4/3\).

\subsection*{R3. Montrer où l'uniformité en coefficients intervient}
\textbf{Traitement : tableaux de paramètres et ordre des choix dans T1.}
Les preuves proches de T1.4 et T1.7 distinguent dimension ambiante,
degré, hauteur des points, hauteur des coefficients et normalisation
des boîtes. La dépendance de l'énoncé de Reuss en son paramètre fixé
est signalée ; l'uniformisation utilisée ici est rattachée au calcul
du compagnon. Elle n'est pas présentée comme une conséquence gratuite
d'un théorème cité avec coefficients fixés.

\clearpage
\section*{3. Le registre des marges}
Dans ce tableau, un exposant \(\beta\) signifie une contribution
majorée par \(N^{-\beta+\text{perte}}\). La perte à comparer à la marge
est celle de la \emph{branche complète}, après produits, propagations
et choix de configurations.

\small
\begin{tabularx}{\textwidth}{@{}>{\raggedright\arraybackslash}X>{\centering\arraybackslash}p{.19\textwidth}>{\centering\arraybackslash}p{.19\textwidth}@{}}
\toprule
Comparaison & Écart exact & Usage\\
\midrule
Cible non alignée au-dessus d'une paire rationnelle :
\(11/8\) contre \(41/30\)
& \(1/120\) & Triplets\\[5pt]
Queue uniforme à faible nullité :
\(27/20\) contre \(4/3\)
& \(1/60\) & Reconstruction\\[5pt]
Activité des triplets :
\(41/30\) contre \(4/3\)
& \(1/30\) & Majorant scalaire\\[5pt]
Triplets incompatibles :
\(3/2\) contre \(41/30\)
& \(2/15\) & Signe du cumulant\\
\bottomrule
\end{tabularx}
\normalsize

Les quatre exposants de décroissance issus des certificats uniformes
après paiement du facteur de rang sont
\[
 \frac32,\qquad \frac{1106}{375},\qquad
 \frac{27}{20},\qquad \frac{2279}{1500}.
\]
Leurs gains sur \(4/3\) sont respectivement
\[
 \frac16,\qquad \frac{202}{125},\qquad
 \frac1{60},\qquad \frac{93}{500}.
\]
Les produits de deux activités de paires et les primitives à quatre
blocs sont en revanche déjà à l'échelle \(N^{-4/3+o(1)}\).
Ils sont conservés dans le reste cible ; aucune marge strictement
positive ne leur est attribuée.

\subsection*{Ordre des choix}
On prescrit d'abord \(C\) et la précision finale \(\eps_0>0\).
On choisit ensuite une précision auxiliaire
\[
 0<\tau<\min(\eps_0/4,\,1/960).
\]
Chaque estimation auxiliaire est demandée avec une perte totale au
plus \(\tau\), en répartissant cette précision entre ses propres
paramètres internes. Une multiplication de deux estimations peut
donc consommer \(2\tau\). Les seuils géométriques et de rang ne sont
pas tous identifiés à ce même nombre.

Pour la borne inférieure, l'ordre de Hölder \(k\) est fixé après
\(\eps_0\), avec \(4/[3(k-1)]\le\eps_0\), puis \(N\) tend vers
l'infini. Pour les activités connexes de la reconstruction, les ordres
jusqu'à 256 restent séparés d'une queue uniforme ; 256 est un seuil
suffisant, sans revendication d'optimalité.
Le fichier \texttt{editorial/reviews/margin\_budget.json} contient
les calculs rationnels et leur emplacement dans les preuves.

\clearpage
\section*{4. Dépendances, portée et reproductibilité}

\subsection*{R4. Identifier les chemins de preuve}
\textbf{Traitement : carte à flèches en section 1.5 et guide du compagnon.}
Les moments fixés de \(Z_N\) sont prouvés avant l'approximation en loi.
Les moments des sélections courtes, qui portent sur des valeurs
polynomiales de la fonction multiplicative, sont un autre outil,
avec leur propre plage de degrés.

Le contrôle des triplets alimente la reconstruction. Le lemme quartique
utilise les estimations de paires, de triplets et de déficits d'extension,
puis le minorant utilise les moments fixés indépendants. Il n'emprunte
pas la conclusion de reconstruction pour établir ces moments.
Les erreurs scalaires pondérées sont, elles, des conséquences du
théorème principal et de ces mêmes moments.

\subsection*{R5. Maintenir les limites de chaque optimalité}
\textbf{Traitement : conservé et vérifié dans l'introduction et la discussion.}
Le taux \(4/3\) est une optimalité de puissance pour la référence
quadratique exacte. Il ne fournit ni constante principale ni facteur
logarithmique du reste. La puissance \(41/30\) du troisième cumulant
ne constitue pas davantage un équivalent complet.

La référence de Hermite exacte contient déjà
\(\alpha_N^2\Pi_{\mu_N}*\Delta^{*4}/8\). Le test quartique utilisé
pour \(Q_{\mu_N,\alpha_N}\) ne donne donc pas le même minorant
pour Hermite. Les théorèmes de canaux conservent \(h\) fixé.
Le résultat positionnel garde les positions entières non signées
avec erreur de puissance \(4/5\), sans signes ni longueurs excédentaires.

\subsection*{R6. Classer les prolongements et donner une commande unique}
\textbf{Traitement : tableau dans le guide du compagnon et vérification regroupée.}
La preuve positionnelle de T19 est nécessaire au lemme 20.1 de l'article.
Les autres résultats de T19 et les sections T20--T23 apportent des
prolongements ou explicitent des limites ; ils ne sont pas des entrées
du théorème scalaire optimal. Cette classification décrit leur rôle
logique au sein du même travail.

Depuis la racine des sources, une seule commande lance les compilations,
les cinq diagnostics finis et le contrôle des références :
\begin{center}
\texttt{python3 check\_all.py}
\end{center}
Le résultat structuré est écrit dans
\texttt{editorial/verification\_run.json}. Les journaux de compilation
restent dans le répertoire de travail \texttt{build/}.
Le README détaille les dépendances et les identités examinées.
Ce dispositif contrôle les documents et certains calculs finis ; il
n'est pas un vérificateur des démonstrations asymptotiques.

\clearpage
\section*{5. Réponse à l'audit bibliographique}
Les six ajouts sont numérotés 26--31 après les 25 références déjà
présentes, dont le contenu et l'ordre sont conservés. Les numéros
ci-dessous désignent la bibliographie commune aux deux manuscrits.

\small
\begin{tabularx}{\textwidth}{@{}>{\raggedright\arraybackslash}p{.28\textwidth}>{\raggedright\arraybackslash}X@{}}
\toprule
Source ajoutée & Attribution et emplacement\\
\midrule
{[26]} Griesmer (1960)
& Borne binaire et preuve par code résiduel, avant le lemme 2.1.
Le support arithmétique minimal et l'application aux fenêtres restent
explicitement distingués de l'inégalité classique.\\[6pt]
{[27]} Zeilberger (1997)
& Identité d'arrêt en section 5. La correspondance avec la lace-map,
ses points fixes et ses ensembles compatibles est donnée après le
lemme 5.1 ; l'identité générale n'est pas revendiquée comme nouvelle.\\[6pt]
{[28]} Cilleruelo--Garaev (2011)
& Comptage quadratique uniforme de T1.2, avec référence au lemme 4
et à la proposition 1 de la version précisée. Le cas non carré et
le cas traité par factorisation sont séparés.\\[6pt]
{[29]} Tao (2026, v3)
& Formulation complémentaire au lemme 2.10, citée dans T1.2.
Cette référence ne remplace pas l'antécédent de Cilleruelo--Garaev,
que le texte de Tao cite lui-même.\\[6pt]
{[30]} Čekanavičius--Vellaisamy (2015)
& Comparaison des corrections quadratique et exponentielle :
introduction, ouverture du quartique et discussion. Leurs hypothèses
de dépendance locale ne sont pas transposées au modèle de P1.\\[6pt]
{[31]} Chhaibi et coauteurs (2020)
& Contexte des équivalents de distance en convergence mod-\(\phi\).
Le paramètre divergent et la régularité des résidus de Fourier du
théorème 3.11 sont distingués du régime de moyennes bornées de P1.\\
\bottomrule
\end{tabularx}
\normalsize

\subsection*{Ce qui a été vérifié dans les sources}
Les localisateurs ont été confrontés aux versions primaires indiquées :
§2 et (2.10) de Griesmer ; formule (*) de Zeilberger ;
§3 de Cilleruelo--Garaev ; lemme 2.10 de Tao v3 ;
théorèmes 3.1--3.2, (3.5)--(3.6), du journal ALEA ;
théorème 3.11 et hypothèse H2 du texte publié des Annales de l'Institut Fourier.

L'article original de Griesmer a été lu sur son scan, bien que son
résolveur DOI n'ait pas été accessible pendant ce contrôle.
Les versions, URL, localisateurs et limites sont consignés dans
\texttt{editorial/reviews/bibliography\_verification.json}.
Il s'agit d'une vérification ciblée des ajouts, sans prétention à une
recherche exhaustive de priorité.

Le crédit donné à \emph{Paper C v3} est conservé : la loi de Poisson
du champ marqué et l'apparition de l'interaction \(2:3\) ne sont pas
présentées comme des nouveautés de P1.

\clearpage
\section*{6. Contrôles, sources et circulation}

\subsection*{Intégrité des deux manuscrits}
La table de correspondance V1.4--V1.5 suit les
\PreservedLabelCount{} anciens repères et les
\NumberedResultCount{} énoncés numérotés. Le contrôle vérifie leur
conservation, l'absence de doublons, la stabilité des numéros des
énoncés, la résolution des renvois et les cibles des liens PDF.
Les compilations finales contrôlent \CrossDocumentLinkCount{} liens
entre documents et \InternalLinkCount{} liens internes.

Les deux documents scientifiques doivent être conservés ensemble
avec leurs noms fournis pour utiliser les liens croisés. Les références
imprimées, notamment les numéros commençant par T, restent utilisables
sans cette fonction du lecteur PDF.

\subsection*{Les cinq diagnostics finis}
\begin{enumerate}
\item \textbf{Marginales hétérogènes :} huit lois jointes finies de Bernoulli,
112 contrôles exacts des probabilités, coefficients centrés, référence
de Bernoulli avec covariances et identités de cumulants factoriels jusqu’à l’ordre quatre.
\item \textbf{Arrêt héréditaire :} toutes les classes héréditaires
contenant l'ensemble vide sur au plus quatre sites et tous leurs
états d'indicateurs ; exemple à suffixe dépendant, identité de
paire retenue et seuils finis de positivité.
\item \textbf{Coefficients :} identités de canaux de hauteurs
2 à 12, longueurs jusqu'à 96, deux conventions d'extrémité ;
placements signés, normes de la deuxième différence de Poisson aux moyennes 1 et 2, facteur
\(1/(12e)\) à moyenne unitaire.
\item \textbf{Canaux et triplets :} identités de hauteur au plus 10,
préfacteurs rationnels des constantes de phase aux hauteurs 4, 5 et 6,
triplets primitifs dyadiques
\(2\le u<150\), \(u<v<w<2u\), et nullités du témoin
\(10:12:15\) pour les longueurs 31 à 300.
\item \textbf{Construction biquadratique :} une instance symbolique
exacte de l'auxiliaire, avec calculs rationnels et vérifications
de ses annulations.
\end{enumerate}
Les cinq scripts sont relancés par la commande unique. Leur domaine
fini est explicite. En particulier, l'énumération des triplets ne
remplace pas la preuve générale de la borne \(19/30\), et un
certificat symbolique isolé ne prouve pas l'uniformité en degré.

\subsection*{Contenu de l'archive}
L'archive contient les sources des trois PDF, la bibliographie commune,
les commandes de compilation et de contrôle, les diagnostics et leurs
résultats, les correspondances V1.3--V1.4 et V1.4--V1.5, les comptes
rendus ciblés et un manifeste SHA-256. Les anciens comptes rendus
de V1.4 sont identifiés comme historiques.

Les relectures assistées et les calculs exacts renforcent la traçabilité
de cette révision. Ils ne constituent ni une expertise humaine
indépendante ni une certification par assistant de preuve.
L'auteur reste responsable des résultats et de toute soumission.
\end{document}
